\section{Discussion}
\label{sec:discussion}

\subsection{What the Results Mean for a Self-Test}

The results support a self-test in two levels for the realistic circuit. The first level, the
compliance decision, can be set to let through about \SI{1}{\percent} of the non-compliant
circuits while rejecting \SI{8}{\percent} of the compliant ones, or \SI{2}{\percent} when porous
dry electrodes are excluded. A limit test on the same measurements cannot approach this, because
it asks whether the circuit looks like a healthy one, which is neither necessary nor sufficient
for compliance. The second level should report an ambiguity group, not a component. The groups
are known from the nominal circuit before any fault is simulated, group-level location holds for
fault magnitudes the model has not seen, and for repair the difference is small: in \cA\ a group
names two to four parts that set the same quantity.

The comparison of labels has a direct practical reading. A detector trained in the way that is
usual in the diagnosis literature, on whether a component deviates, would take out of service
more than a third of the instruments that still meet the standard. This is not a weakness of the
classifier but of the question it is asked; the same model, asked about compliance, is right.

The cost of the self-test is low. Three measurements and three seconds give the decision, and the
most expensive ones, the tones below the high-pass corner, are not needed. This matters in
practice because the self-test has to run with the patient connected and without a signal.

\subsection{The Electrode as a Confounder}

The electrode limits both levels. A healthy circuit on a high-impedance electrode measures a
reduced gain through the calibration source, which is what a gain fault looks like. On \cA\ this
accounts for eight of the nine invisible fault conditions, and false rejects at the
\SI{1}{\percent} escape target fall from 7.8 to \SI{1.9}{\percent} when porous electrodes are
left out. And a
degraded contact can only be recognized against the normal contact of the same subject, which a
single self-test does not know. Two consequences follow. The lead-off tone is among the first
measurements selected, plausibly because it tells the model how much of the attenuation the
electrode explains; that reading of the selection is ours and was not tested separately. And a self-test that kept a reference per patient, or that could disconnect the
electrodes from the calibration path, would remove most of the difficulty; neither option was
studied here.

\subsection{Realistic and Benchmark-Like Circuits}

The two circuits share the signal chain and differ in how the instrumentation amplifier is built,
and that difference changes the conclusions. The discrete amplifier adds seven matched resistors
that all act on the gain, which merge into one group of nine with the gain stage, and its
common-mode sensing makes several faults of the DRL stage both non-compliant and invisible. The
integrated amplifier removes the first problem and, in this design, the second. A method
evaluated only on the benchmark-like circuit would have been judged unable to reach a low escape
rate; evaluated on the realistic one it can. This supports the argument of~\cite{chen2025} that
the choice of circuit conditions the conclusions, with a different remedy: instead of redesigning
the circuit to be separable, the ambiguity that the given circuit has is computed and the
diagnosis is reported at that level.

The invisible DRL faults also show a limit of the measurements. Raising the amplitude or the
duration of the common-mode tone does not help, whereas the dc reading of the DRL output is the
third measurement selected for the decision on \cB\ in every repetition. Where a spare converter
channel is available, it is better spent there.

\subsection{Limitations}

The study is based on simulation only, with behavioral models of the amplifiers and ideal
self-test sources; how well models trained on simulated data transfer to hardware has not been
tested. Faults are single, and aging is not modeled. The frequency-domain features are
small-signal values. The dry-electrode model rests on six measured subjects per material; the
gel electrode has no spread between subjects, and the spread between the electrodes of one case
is assumed. The noise specification is computed from the rms value and the input range from the
operating point, instead of from the time-domain tests of the standard. Finally, the
specifications are those of performance: some hard faults leave the circuit compliant while
removing a protection (a short across the resistor that limits the current into the right-leg
electrode, for instance), and a self-test aimed at safety would have to label them separately.
